\begin{helper}[Exit localization for the two order-nine blocks]
Let $(T,e_0,\varphi,N)$ be a balanced normal globally maximum triple with
$T.\mathsf{length}=8$ and $N.\mathsf{height}=5$.  Let
$\gamma_0,\gamma_1$ be distinct colors such that the identities in
positions $4,5$ have color $\gamma_0$, the identities in positions $6,7$
have color $\gamma_1$, each $\gamma_i$ is missing at one of the first five
vertices, and the colored base identities in positions $1,2,3$ avoid both
$\gamma_i$.  Then every Scheide exit vertex of $T$ misses one of
$\gamma_0,\gamma_1$.
\end{helper}
\begin{proof}
Write $y_r=T.\mathsf{vertex}(r)$.  The normal height-five path in
\noderef{NormalGlobalMaximumTriple} has first normal color $\alpha$
missing at $y_0$ and second normal color $\beta$ missing at $y_1$.
Its exact support and endpoint fields, together with properness, show
that the colored base identities in positions $1,2,3$ alternate between
$\alpha$ and $\beta$.  The base-avoidance hypothesis therefore makes
both $\alpha$ and $\beta$ distinct from both block colors; injectivity of
the block-color map makes $\gamma_0\ne\gamma_1$.

The first identity is $e_0$ by \noderef{TreeSequence} and is the unique
uncolored identity by the $k$-triple inside the globally maximum field.
The preceding base-color calculation and the two block-color hypotheses
therefore give the exact equality
\[
 \noderef{TreeUsedColors}(\varphi,T)
   =\{\alpha,\beta,\gamma_0,\gamma_1\}.
\tag{1}
\]
Choose $z_i$ among $y_0,\ldots,y_4$ with $\gamma_i$ missing at $z_i$,
as supplied by the hypotheses.  The elementary clause of
\noderef{ScheideGlobalMaximumOrder}, expanded through
\noderef{ColoredVertexSetElementary}, says that missing-color sets at
distinct displayed vertices are disjoint.  Consequently $y_0$ is the
only displayed vertex that can miss $\alpha$, $y_1$ the only one that can
miss $\beta$, and $z_i$ the only one that can miss $\gamma_i$.

We next count missing colors.  The degree bound in the $k$-triple and
\noderef{MaximumDegree} give
$d_G(v)\leq\Delta(G)\leq k-1$ for every vertex $v$, with degree interpreted
by \noderef{Degree}.  Properness makes the colors of the colored
identities incident with $v$ pairwise distinct.  If $v$ is not an end of
$e_0$, every incident identity is colored, so $v$ misses at least
$k-d_G(v)\geq1$ color.  If $v$ is an end of $e_0$, exactly that incident
identity is uncolored and every other incident identity is colored, so
$v$ misses at least $k-d_G(v)+1\geq2$ colors.  The attachment equation at
edge index zero and the first-edge equation in
\noderef{TreeSequence} identify the two ends of $e_0$ as $y_0,y_1$.

Let $v$ be a displayed vertex different from both $z_0,z_1$.  If $v$ is
not an end of $e_0$, then it is also different from $y_0,y_1$.  Its
guaranteed missing color is therefore none of the four colors in (1), by
elementarity and the four displayed witnesses.  If $v$ is an end of
$e_0$, it is $y_0$ or $y_1$.  At most its corresponding one of
$\alpha,\beta$ can be missing among the four colors in (1), because $v$
is different from both $z_i$; the lower bound of two missing colors again
supplies a missing color outside (1).  In either case $v$ misses a color
not used on any colored tree identity, hence a color in
\noderef{TreeFreeColors}.

By \noderef{ScheideExitNoFree}, a Scheide exit vertex has its missing
colors disjoint from the tree-free colors.  Therefore every exit is one
of $z_0,z_1$.  At $z_i$ the color $\gamma_i$ is missing by its defining
choice, which proves the claimed exit coverage.
\end{proof}
